\chapter{АНАЛИЗ ПРЕДМЕТНОЙ ОБЛАСТИ И ПОСТАНОВКА ЗАДАЧИ}

% ─────────────────────────────────────────────────────────────────
\section{Проблема и её масштаб}
% ─────────────────────────────────────────────────────────────────

\subsection{Болевые точки авторов курсов}

Разработка корпоративного учебного курса остаётся трудоёмким и слабо
масштабируемым процессом. Даже при наличии готового экспертного материала
необходимо определить целевую аудиторию, сформулировать учебные цели,
разбить содержание на модули и уроки, подобрать формат подачи,
подготовить оценочные материалы и синхронизировать всё это
с~требованиями LMS.
По~оценке Chapman Alliance, производство одного часа eLearning-контента
требует от~42 до~490 часов работы в~зависимости от~сложности сценария,
мультимедийности и глубины методической проработки; для большинства
практических проектов реалистичный коридор составляет 40--200 часов
на~один час готового обучения~\cite{chapman_elearning_hours_2010}.

Для российского B2B-сегмента проблема усугубляется дополнительно.

\begin{itemize}
  \item \textbf{Дефицит методических компетенций.} Эксперт предметной
    области знает материал, но не всегда умеет переводить знание
    в~педагогически корректную структуру курса. Другой проблемой становится нехватка или перегрузка специалистов
    в~области методики преподавания. У компаний зачастую нет финансовой возможности иметь собственный методический отдел.
  \item \textbf{Высокая стоимость внешней разработки.} На рынке
    корпоративного обучения разработка одного курса силами подрядчика
    обычно оценивается в диапазоне 100--300~тыс.\ руб.\ за проект,
    что делает частое обновление контента экономически невыгодным
    для малого и среднего бизнеса~\cite{smart_ranking_elearning_2024}.
  \item \textbf{Быстрое устаревание источников.} Регламенты, продуктовые
    описания и внутренние правила обновляются быстрее, чем команда
    успевает вручную пересобирать учебные материалы, из-за чего курс
    быстро расходится с~актуальной версией исходного документа.
\end{itemize}

\subsection{Разрыв между экспертизой и образовательной упаковкой}

В корпоративной среде инициатором
курса обычно выступает владелец продукта, руководитель функции, HR или
L\&D-менеджер. Первый хорошо знает предметную область, второй понимает
бизнес-контекст, однако ни один из них не обязательно владеет
инструментами педагогического дизайна на уровне, достаточном для создания
масштабируемого цифрового курса.

На практике это проявляется так: черновой
материал остаётся на уровне «экспертного конспекта» без явных
учебных результатов и уровней освоения. Структура курса
строится от имеющихся документов, а не от целей обучения, что снижает
практическую ценность. Сопровождение же требует
подключения методиста, зачастую не входящего в команду.

\subsection{Подтверждение проблемы: результаты исследования клиентов}

Для первичной проверки гипотез о проблеме проведено 20 глубинных
B2B-интервью с~HR-специалистами, L\&D-менеджерами, методистами и
корпоративными тренерами по~подходу Mom~Test~\cite{fitzpatrick_mom_test_2013}.
Руководство интервью, логика кодирования ответов и шаблон таблицы выводов
вынесены в приложение~А.

Количественные результаты:

\begin{itemize}
  \item \textbf{70\,\%} респондентов называют время подготовки курса
    главным ограничением запуска нового обучения;
  \item \textbf{50\,\%} респондентов уже привлекают внешних подрядчиков
    или отдельных методистов для упаковки материала;
  \item \textbf{60\,\%} готовы рассмотреть ИИ-инструмент как часть
    текущего процесса разработки, если он не ухудшает качество
    структуры и позволяет сохранить контроль над итоговым результатом.
\end{itemize}

Распределение ответов респондентов представлено на рисунке~\ref{fig:custdev-bars}.

\begin{figure}[htbp]
\centering
\begin{tikzpicture}
\begin{axis}[
    xbar,
    width=0.85\textwidth,
    height=5cm,
    xlabel={Доля респондентов, \%},
    xmin=0, xmax=100,
    ytick=data,
    yticklabels={
        {Готовы рассмотреть ИИ-инструмент},
        {Привлекают внешних подрядчиков},
        {Время разработки --- главное ограничение}
    },
    y tick label style={font=\small, anchor=east},
    nodes near coords,
    nodes near coords align={horizontal},
    nodes near coords style={font=\small, append after command={node[anchor=west]{\,\%}}},
    bar width=14pt,
    enlarge y limits=0.35,
    every axis plot/.append style={fill=teal!80!black, draw=teal!60!black},
]
\addplot coordinates {(60,0) (50,1) (70,2)};
\end{axis}
\end{tikzpicture}
\caption{Результаты Customer Development (n=20 B2B-интервью)}
\label{fig:custdev-bars}
\end{figure}


Качественный вывод: обучающиеся находятся в
профиците информации и дефиците времени. Отсюда запрос --- быстро превратить регламент,
брендбук или продуктовый бриф в обучение с понятной логикой, проверяемыми
учебными целями и практическим результатом.

\subsection[Решенческие интервью и критерии решения]{Решенческие интервью и критерии соответствия решения проблеме}

Серия решенческих интервью проведена на базе функционального прототипа ИИ-методиста
(сценарий демонстрации --- приложение~Б). Проверяемые
критерии соответствия решения проблеме:

\begin{enumerate}
  \item готов ли руководитель L\&D-отдела встроить систему
    в~текущую ручную стадию подготовки черновой структуры курса
    как~инструмент ускорения и~снятия рутины;
  \item готов ли методист использовать систему на реальном документе
    своей организации;
  \item экономически эффективно ли решение по отношению к
    стоимости аутсорс-разработки или внутреннего
    методического ресурса.
\end{enumerate}

\subsection{Рыночный контекст: TAM / SAM / SOM}

Глобальный рынок онлайн-обучения растёт:
по оценке Grand View Research, его объём составил
\$299{.}67~млрд в~2024~году и может достичь \$842{.}64~млрд к~2030~году
при CAGR около 19\,\%~\cite{grandview_elearning_2025}. Это задаёт
широкий TAM для решений, ускоряющих подготовку цифровых
образовательных материалов.

Для оценки сегмента в~России в работе используется
более консервативная логика. По данным Smart Ranking, объём
российского EdTech-рынка в~2023~году достиг 119{.}33~млрд~руб.,
а доля сегмента «разработчики и платформы» составляла 13\,\%
от общей структуры рынка~\cite{smart_ranking_elearning_2024}.
Актуальные обзоры СберОбразования за 2025~год также фиксируют
рост внимания к ИИ-инструментам, персонализации и автоматизации
образовательных процессов~\cite{sber_education_trends_q2_2025,sber_education_trends_q3_2025}.
Этот слой ближе всего по природе
к ИИ-методисту; рынок инструментов поддержки
разработки курсов оценивается примерно в~15{.}5~млрд~руб.

При целевой доле 0{.}1\,\% от~SAM в~горизонте трёх лет
SOM составляет около 15{.}5~млн~руб.\ в год. При базовой B2B-подписке порядка
15~тыс.\ руб.\ в месяц это соответствует примерно
85--90 активным корпоративным клиентам. Такой масштаб достижим
для нишевого B2B-продукта.

Структура рыночной оценки показана на рисунке~\ref{fig:tam-sam-som}.

\begin{figure}[htbp]
\centering
\begin{tikzpicture}
% TAM — widest
\fill[teal!70!black, rounded corners=2pt] (-6,0) rectangle (6,1.2);
\node[font=\bfseries\small, text=white] at (0,0.6) {TAM (глобальный): \$842{,}64 млрд к~2030 г.};

% SAM — medium
\fill[teal!50!black, rounded corners=2pt] (-4,-1.8) rectangle (4,-0.6);
\node[font=\bfseries\small, text=white] at (0,-1.2) {SAM (EdTech-инструменты, РФ): 15{,}5 млрд руб.};

% SOM — narrowest
\fill[teal!30!black, rounded corners=2pt] (-2.2,-3.6) rectangle (2.2,-2.4);
\node[font=\bfseries\small, text=white] at (0,-3.0) {SOM (3 года): 15{,}5 млн руб.};

% Connecting dashed lines (funnel edges)
\draw[gray, dashed, thick] (-6,0) -- (-4,-0.6);
\draw[gray, dashed, thick] (6,0) -- (4,-0.6);
\draw[gray, dashed, thick] (-4,-1.8) -- (-2.2,-2.4);
\draw[gray, dashed, thick] (4,-1.8) -- (2.2,-2.4);
\end{tikzpicture}
\caption{Оценка рынка: TAM / SAM / SOM}
\label{fig:tam-sam-som}
\end{figure}


% ─────────────────────────────────────────────────────────────────
\section{Анализ существующих продуктов}
% ─────────────────────────────────────────────────────────────────

\subsection{Прямые конкуренты}

\textbf{Coursebox~AI} позиционируется как конструктор курсов на~основе ИИ и
позволяет преобразовывать документы, слайды, видео и URL в курс,
автоматически генерировать тесты, видео, изображения и публиковать
результат в собственную LMS либо экспортировать в SCORM/LTI
форматах~\cite{coursebox_ai_2024}. Это один из наиболее близких
по~ценностному предложению продукт на~рынке.
Однако его логика остаётся в~первую очередь ориентированной
на~быструю упаковку содержимого:
сильная сторона продукта --- скорость, а не явная
педагогическая методология или контроль корректности фактов
по отношению к исходному документу.

\textbf{Articulate 360 AI} добавляет ИИ-функции в зрелую среду
корпоративного контента: создание структуры курса, чернового текста,
перевода, изображений и озвучки внутри экосистемы Articulate.
Продукт хорошо встроен в привычный рабочий процесс методиста,
но остаётся инструментом повышения производительности разработчика курса,
а не самостоятельным методическим конвейером с RAG-слоем,
проверкой фактов и локальным развёртыванием~\cite{articulate_ai_2025}.

\textbf{iSpring Suite AI} развивает похожую траекторию:
в~2025~году в продукт добавлены AI Course Creation Helper,
AI Translator, AI Image Generator, а также функции синтеза речи и
другие инструменты ускорения разработки курса~\cite{ispring_suite_ai_2026}.
Сильной стороной iSpring является зрелая экосистема для корпоративного
обучения и русскоязычный рынок. Ограничение состоит в том, что ИИ
остаётся надстройкой над средой разработки контента, а не автономным
мультиагентным конвейером, принимающим на вход первичный документ
и проводящим его через полную цепочку анализа, проектирования и
обновления курса.

\textbf{Synthesia} лидирует в классе ИИ-видеогенерации для обучения
и корпоративных коммуникаций: её ценность сосредоточена в создании
аватарных роликов, озвучки и мультиязычного видео-контента
для обучения и L\&D-задач~\cite{synthesia_ai_training_2025}.
Но Synthesia не закрывает задачу проектирования учебной структуры.

\subsection{Смежные решения}

\textbf{Google NotebookLM} --- сильный смежный продукт для работы
с источниками, который умеет строить краткие выжимки, аудиообзоры
и видеообзоры по загруженным материалам. В~2025~году Google
расширил генерацию аудио и видео обзоров до 80+ языков, сделав
их заметно реалистичнее и полезнее для изучения сложных источников
~\cite{notebooklm_audio_2025}. Однако NotebookLM не предназначен
для создания LMS-совместимой структуры курса, педагогического
дизайна и контроля соответствия учебных целей.

\textbf{GetCourse} является одним из крупнейших российских игроков
рынка LMS/онлайн-школ: платформа сообщает о 50{,}000+ клиентах,
16~млн учеников и совокупном обороте школ на платформе в
158~млрд~руб.\ за 2025~год~\cite{getcourse_platform_2026}.
Для настоящей работы GetCourse важен как ориентир класса систем, в которые ИИ-методист потенциально
может встраиваться как сервис генерации и обновления курса.

\subsection{Сравнительная матрица}

В таблице~\ref{tab:competitors} приведено сравнение решений по критериям,
которые оказались значимыми в исследовании клиентов: способность работать
от исходного документа, наличие встроенной педагогической логики,
поддержка русскоязычного контента, суверенитет данных, автоматическое
обновление содержания и интеграция с LMS-средой.

\begin{table}[h!]
  \caption{Сравнительный анализ конкурентных решений}
  \label{tab:competitors}
  \centering
  \small
  \begin{tabular}{|p{3.2cm}|c|c|c|c|c|}
  \hline
  \textbf{Критерий}
    & \textbf{Coursebox}
    & \textbf{Articulate~AI}
    & \textbf{iSpring~AI}
    & \textbf{Synthesia}
    & \textbf{ИИ-методист} \\
  \hline
  Генерация курса из документа
    & +
    & $\pm$
    & $\pm$
    & ---
    & \textbf{+} \\
  \hline
  Явная методология (ADDIE / SAM / BD)
    & ---
    & ---
    & ---
    & ---
    & \textbf{+} \\
  \hline
  Русскоязычный контент
    & $\pm$
    & $\pm$
    & +
    & $\pm$
    & \textbf{+} \\
  \hline
  Локальное развёртывание и~суверенитет данных
    & ---
    & ---
    & ---
    & ---
    & \textbf{+} \\
  \hline
  ИИ-аудио / видео
    & +
    & +
    & +
    & \textbf{+}
    & + \\
  \hline
  Проверка фактов по источнику
    & ---
    & ---
    & ---
    & ---
    & \textbf{+} \\
  \hline
  Автоматическое обновление курса
    & $\pm$
    & ---
    & ---
    & ---
    & \textbf{+} \\
  \hline
  Ревью с~участием человека
    & +
    & +
    & +
    & +
    & \textbf{+} \\
  \hline
  Интеграция с LMS-потоком разработки
    & +
    & +
    & +
    & ---
    & \textbf{+} \\
  \hline
  \end{tabular}
\end{table}

Визуализация сравнения в форме радарной диаграммы приведена на рисунке~\ref{fig:competitor-radar}.

\begin{figure}[htbp]
\centering
\begin{tikzpicture}
\begin{polaraxis}[
    name=radar,
    width=0.78\textwidth,
    height=0.78\textwidth,
    xtick={0,60,120,180,240,300},
    xticklabels={
        {Генерация\\из документа},
        Методология,
        {Русский\\язык},
        {Локальное\\развёртывание},
        {Проверка\\фактов},
        {Автоматическое\\обновление}
    },
    xticklabel style={font=\footnotesize, anchor=center, align=center, inner sep=8pt},
    ymin=0, ymax=1.2,
    ytick={0.25,0.5,0.75,1},
    yticklabels={},
    y grid style={gray!30},
    x grid style={gray!50, thick},
    grid=both,
    legend style={
        at={(0.5,-0.1)},
        anchor=north,
        font=\footnotesize,
        cells={anchor=west},
        legend columns=3,
        column sep=6pt,
    },
]

% ИИ-методист (рисуем первым — фон)
\addplot[teal!70!black, very thick, mark=*, mark size=2.5pt, fill=teal!12!white]
    coordinates {(0,1)(60,1)(120,1)(180,1)(240,1)(300,1)(360,1)};
\addlegendentry{ИИ-методист}

% Coursebox
\addplot[blue!70!black, thick, mark=square*, mark size=2.5pt]
    coordinates {(0,1)(60,0)(120,0.5)(180,0)(240,0)(300,0.5)(360,1)};
\addlegendentry{Coursebox}

% Articulate AI
\addplot[red!70!black, thick, mark=triangle*, mark size=2.5pt]
    coordinates {(0,0.5)(60,0)(120,0.5)(180,0)(240,0)(300,0)(360,0.5)};
\addlegendentry{Articulate AI}

% iSpring AI
\addplot[orange!80!black, thick, mark=diamond*, mark size=2.5pt]
    coordinates {(0,0.5)(60,0)(120,1)(180,0)(240,0)(300,0)(360,0.5)};
\addlegendentry{iSpring AI}

% Synthesia
\addplot[purple!70!black, thick, mark=pentagon*, mark size=2.5pt]
    coordinates {(0,0)(60,0)(120,0.5)(180,0)(240,0)(300,0)(360,0)};
\addlegendentry{Synthesia}

% Подписи шкалы вдоль одного луча (90° = вверх-вправо, ось "Методология")
\node[font=\tiny, fill=white, inner sep=1pt, text=gray!70!black] at (axis cs:30,0.25) {0{,}25};
\node[font=\tiny, fill=white, inner sep=1pt, text=gray!70!black] at (axis cs:30,0.5) {0{,}5};
\node[font=\tiny, fill=white, inner sep=1pt, text=gray!70!black] at (axis cs:30,0.75) {0{,}75};
\node[font=\tiny, fill=white, inner sep=1pt, text=gray!70!black] at (axis cs:30,1.0) {1{,}0};

\end{polaraxis}
\end{tikzpicture}
\caption{Радарная диаграмма сравнения решений}
\label{fig:competitor-radar}
\end{figure}


\subsection{Пробел между проблемой и существующими решениями}

Рынок предлагает и стартапы, и зрелые продукты
для ускорения создания обучения. Однако сохраняется разрыв между генерацией учебных материалов и адаптацией к контексту каждого обучающегося.

Ни одно из рассмотренных решений не закрывает одновременно четыре
требования корпоративного B2B-заказчика:

\begin{enumerate}
  \item работу от внутреннего документа организации как от первичного
    источника;
  \item поддерживать методическую логику проектирования курса для достижения целей обучения;
  \item возможность локального развёртывания и контроля данных;
  \item механизм повторной генерации и обновления курса при изменении
    исходного источника.
\end{enumerate}

Позиционирование решений на карте приведено на рисунке~\ref{fig:positioning-map}.

\begin{figure}[htbp]
\centering
\begin{tikzpicture}
\begin{axis}[
    width=0.85\textwidth,
    height=0.7\textwidth,
    xlabel={Методическая глубина},
    ylabel={Суверенитет данных},
    xmin=0, xmax=10,
    ymin=0, ymax=10,
    grid=major,
    grid style={gray!30},
    xlabel style={font=\small},
    ylabel style={font=\small},
    tick label style={font=\small},
    clip=false,
]

% Regular competitors
\addplot[only marks, mark=*, mark size=3pt, teal!60!black]
    coordinates {(3,1)};
\node[anchor=south west, font=\scriptsize] at (axis cs:3.2,1.1) {Coursebox AI};

\addplot[only marks, mark=*, mark size=3pt, teal!60!black]
    coordinates {(4,1)};
\node[anchor=south west, font=\scriptsize] at (axis cs:4.2,1.1) {Articulate AI};

\addplot[only marks, mark=*, mark size=3pt, teal!60!black]
    coordinates {(4,2)};
\node[anchor=south west, font=\scriptsize] at (axis cs:4.2,2.1) {iSpring AI};

\addplot[only marks, mark=*, mark size=3pt, teal!60!black]
    coordinates {(1,1)};
\node[anchor=south west, font=\scriptsize] at (axis cs:1.2,1.1) {Synthesia};

\addplot[only marks, mark=*, mark size=3pt, teal!60!black]
    coordinates {(2,3)};
\node[anchor=south west, font=\scriptsize] at (axis cs:2.2,3.1) {NotebookLM};

% ИИ-методист — highlighted
\addplot[only marks, mark=*, mark size=5pt, fill=teal!80!black, draw=black, thick]
    coordinates {(8,9)};
\node[anchor=south west, font=\scriptsize\bfseries] at (axis cs:8.2,9.1) {ИИ-методист};

\end{axis}
\end{tikzpicture}
\caption{Карта позиционирования конкурентных решений}
\label{fig:positioning-map}
\end{figure}


\subsection{Тяжелокопируемые конкурентные преимущества}

Уникальность предлагаемого решения состоит не в том, что оно использует
LLM для генерации текста, а в соединении методического проектирования,
источниковой доказательности и воспроизводимой проверки в одном
производственном контуре. В отличие от одношаговых запросов,
инструментов ускоренной разработки курсов и исследовательских
блокнотов с опорой на источник,
ИИ-методист не ограничивается созданием черновика: он сохраняет
доказательный контекст, строит структуру курса до генерации уроков,
вводит точки проверки человеком, рассчитывает знак качества и позволяет
повторить парную проверку на том же корпусе.

Ключевое возражение к новизне решения состоит в том, что современный
ChatGPT уже умеет работать с файлами, искать в вебе и в режиме
глубокого исследования (deep research) формировать отчёты
с источниками~\cite{chatgpt_deep_research_2026}.
Это действительно снижает ценность простых одношаговых запросов как
самостоятельной разработки. Поэтому в настоящей работе уникальность
защищается не самим фактом использования LLM, а тем, что ИИ-методист
превращает генерацию курса в контролируемый производственный процесс:
источник становится первым объектом системы, происхождение результата
сохраняется на уровне блоков, знак качества рассчитывается машинно, а
публикация проходит через контрольные точки с участием человека.

В таблице~\ref{tab:uniqueness-defense} сравниваются не маркетинговые
заявления, а свойства, которые можно проверить в артефактах MVP.
Для ChatGPT учитывается текущий класс возможностей OpenAI: поиск и
глубокое исследование дают ссылки на источники, но сам по себе запрос
не создаёт постоянную цепочку подготовки курса и не гарантирует
доменный контроль качества~\cite{chatgpt_deep_research_2026,chatgpt_accuracy_2026}. Для
NotebookLM признаётся сильная работа с цитатами и опорой на источники,
однако продукт остаётся исследовательским блокнотом, а не конвейером
подготовки корпоративного курса~\cite{notebooklm_chat_2026}. Для
Coursebox и iSpring фиксируются сильные функции разработки курса из открытых
материалов~\cite{coursebox_ai_2024,ispring_suite_ai_2026}; отсутствие
отметки означает только то, что в открытом описании продукта не заявлен
такой же проверяемый механизм.

\begin{center}
  \captionof{table}{Защита уникальности ИИ-методиста через сравнение механизмов доверия}
  \label{tab:uniqueness-defense}
  \scriptsize
  \begin{tabular}{|p{2.45cm}|p{2.75cm}|p{2.85cm}|p{2.65cm}|p{3.35cm}|}
  \hline
  \textbf{Критерий}
    & \textbf{ChatGPT + запрос / B2}
    & \textbf{Coursebox / iSpring}
    & \textbf{NotebookLM}
    & \textbf{ИИ-методист B4} \\
  \hline
  Привязка к источнику
    & Зависит от режима и дисциплины запроса; B2 в эксперименте имеет 0\,\% привязки к источнику.
    & Генерируют курс или структуру из материалов, но карта источников на уровне артефактов не заявлена.
    & Ответы опираются на выбранные источники и цитаты.
    & 100\,\% B4-запусков имеют ссылки и связь блока с источником. \\
  \hline
  Происхождение результата
    & История чата не является воспроизводимым журналом курса.
    & В открытом описании не заявлено происхождение каждого фрагмента курса.
    & Сохраняет источники и заметки блокнота, но не цепочку подготовки курса.
    & Сохраняются JSON-файлы со статусом, результатом и происхождением,
      происхождение блоков и исходные доказательства запусков. \\
  \hline
  Методический конвейер
    & Методология остаётся текстом в запросе и легко теряется при повторе.
    & Сильный слой разработки курса: план, тестовые задания, визуальные материалы, перевод, LMS/SCORM.
    & Сильный исследовательский и учебный помощник, но не LMS-конструктор курса.
    & ADDIE реализован как граф: анализ источника, структура, контент, проверка фактов, педагогическая проверка, верификатор. \\
  \hline
  Проверка фактов и знак качества
    & OpenAI рекомендует проверять важные данные; встроенного доменного знака качества для курса нет.
    & В открытых материалах не заявлен воспроизводимый знак качества по источнику.
    & Цитаты помогают проверять ответ, но не дают педагогический контроль качества.
    & 19 из 19 B4-запусков получили статус «пройдено» по знаку качества; F1 проверки фактов = 0.720. \\
  \hline
  Контрольные точки с участием человека
    & Возможны вручную, но не являются частью исполняемого процесса.
    & Редактирование автором есть, но контрольная точка как состояние конвейера не заявлена.
    & Совместное редактирование заметок, без процесса согласования в LMS.
    & Две обязательные точки проверки человеком до публикации и после верификации. \\
  \hline
  Воспроизводимая цепочка артефактов
    & Нельзя честно повторить парное сравнение без внешней дисциплины логирования.
    & Продуктовые материалы описывают ускорение разработки курса, а не экспериментальную цепочку.
    & Есть артефакты блокнота, но нет цикла парной оценки B2/B4.
    & Парный B2/B4 на 19 документах, контроль готовности, исходные доказательства и приложения Ж--К. \\
  \hline
  \end{tabular}
\end{center}

Из сравнения следует, что преимущество ИИ-методиста защищается не
доступом к LLM, а воспроизводимым контуром доверия. ChatGPT может быть
сильным генератором черновика и исследовательским помощником, но
ИИ-методист делает проверяемым весь жизненный цикл учебного артефакта.
В системе явно заданы источник, методическая рамка, цепочка происхождения
результата, автоматическая проверка фактов, знак качества, контрольные
точки с участием человека и парная цепочка B2/B4. Именно эта совокупность даёт долгосрочно развиваемое
преимущество:
каждый новый агент, модель или методология могут подключаться к тем же
контрактам доверия, а качество версии проверяется не впечатлением от
ответа, а повторяемым набором артефактов.
При этом парный B2/B4 на 19 документах доказывает не рыночное
превосходство над аналогами, а работоспособность механизма доверия
в пределах лабораторной проверки MVP. Для сильного сравнительного
утверждения требуется отдельная сравнительная проверка на открытых курсах
$n \geq 30$ с теми же исходными артефактами и явно описанной базовой методикой.

Тяжелокопируемое преимущество формируется на основе четырёх элементов:

\begin{itemize}
  \item \textbf{методологического подхода} --- вариативные педагогические
    рамки ADDIE, SAM и Backward Design становятся параметрами
    агентного конвейера, а не внешним текстовым комментарием;
  \item \textbf{подхода доказательности от источника} --- каждый результат связан
    с источником через ссылки, происхождение блоков, проверку фактов и
    знак качества, что снижает риск неподтверждённой генерации;
  \item \textbf{инфраструктурного подхода} --- быстрый старт в формате
    самостоятельного запуска с~возможностью перехода к~локальному
    развёртыванию на собственной ИИ/LMS-инфраструктуре без вывода данных
    клиента в облако Adapstory;
  \item \textbf{платформного подхода} --- масштабируемой работы человека с ИИ для автоматизации актуализации, проверки фактов, сопровождения
   и других рутинных операций на базе общей истории запусков и
   машиночитаемых доказательных артефактов.
\end{itemize}

% ─────────────────────────────────────────────────────────────────
\section{Обзор научных методов и подходов}
% ─────────────────────────────────────────────────────────────────

\subsection{Формулировка исследовательского вопроса}

Исследовательский вопрос:

\textit{«В какой мере мультиагентная архитектура с сохраняемым состоянием
и вариативной методологической моделью превосходит одношаговый подход при
поддержке методического проектирования корпоративных и учебных материалов
по показателям точности, покрытия уровней таксономии Блума
и времени подготовки первого черновика?»}

Работа находится на стыке трёх линий:
применения LLM в проектировании образовательных материалов,
мультиагентной оркестрации и метрик педагогического качества.

\subsection{LLM для поддержки проектирования образовательных материалов}

Применение LLM в образовании стало особенно заметным после 2022~года.
Kasneci~et~al.~\cite{kasneci_llm_education_2023} показывают, что
большие языковые модели способны существенно ускорять создание
объяснений, учебных заданий и материалов для самостоятельной работы.
Pardos~\cite{pardos_llm_instructional_2023} и Tack~et~al.
~\cite{tack_aied_2022} указывают, что наибольший потенциал LLM
возникает в задачах, где требуется быстро построить первую версию
обучающего материала и затем доработать её с участием человека.

Однако у большинства существующих работ есть повторяющееся ограничение:
генерация строится как одношаговый сценарий или короткая цепочка запросов.
В результате модель выдаёт связный текст, но не обязательно
поддерживает дисциплину педагогического проектирования: учебные цели оказываются
слабо связаны с содержанием, а уровни когнитивной сложности
не распределяются осмысленно.

\subsection{Мультиагентные системы для поддержки проектирования курса}

Мультиагентные LLM-системы декомпозируют сложную задачу
на специализированные роли и состояния.
Wu~et~al.~\cite{wu_autogen_2023} показали, что распределение
ролей между агентами повышает управляемость и качество
LLM-приложений. Shinn~et~al.~\cite{shinn_reflexion_2023} предложили
паттерн Reflexion, а Wang~et~al.~\cite{wang_plan_and_execute_2023}
описали Plan-and-Execute как схему разделения планирования и исполнения.

Проектирование курса ---
\textit{многошаговая} задача: разобрать источник, сформулировать
структуру, произвести контент, проверить соответствие источнику,
при необходимости вернуть в цикл проверка. Это ближе к графу
с~сохраняемым состоянием, чем к диалоговому агенту.

\subsection{RAG в образовании}

RAG-подходы снижают число галлюцинаций за счёт опоры на внешний
корпус знаний~\cite{lewis_rag_2020}. Для учебного курса это критично:
содержание должно воспроизводить смысл исходного документа,
а не просто звучать убедительно.

Es~et~al.~\cite{es_ragas_2023} предложили фреймворк Ragas, который
позволяет измерять достоверность относительно источника
(faithfulness), релевантность ответа (answer relevancy)
и точность контекста (context precision).
Эти метрики хорошо подходят для промежуточной оценки RAG-слоя
ИИ-методиста, так как они фиксируют не только качество ответа,
но и качество связи ответа с источником.

\subsection{Генерация мультимедиа для обучения}

Генерация мультимедийных артефактов ---
иллюстраций, голосовых дорожек, подкастов --- опирается на диффузионные модели
(DALL-E~3, FLUX.1)~\cite{dalle3_2023, flux_2024} и
TTS-решения (Silero, ElevenLabs)~\cite{silero_tts_2024, elevenlabs_2024}.

В настоящей работе мультимедиа --- вспомогательный слой.
Научная новизна не в генерации изображений или аудио как таковой,
а в способе их встраивания в методический конвейер.

\subsection{Методологии проектирования курсов в контексте LLM}

ADDIE~\cite{addie_molenda_2003}, SAM~\cite{allen_sam_2012} и
Backward Design~\cite{wiggins_backward_design_2005} представляют
три различных логики педагогического проектирования.
ADDIE задаёт линейно-этапный подход, SAM делает акцент на быстрых
итерациях, а Backward Design строит курс от измеримых результатов
обучения к структуре и заданиям.

Для LLM-автоматизации каждая из них задаёт различный
способ управления генерацией: что считать первичным --- анализ
аудитории, ожидаемые результаты или итеративное прототипирование.
Методология становится параметром архитектуры, а не внешним комментарием.

\subsection{Оценка качества результатов, подготовленных с~участием ИИ}

Стандартные LLM-метрики слабо подходят для оценки
учебного курса. В работе используется
комбинация трёх классов:

\begin{itemize}
  \item \textbf{структурные метрики}: валидность структуры, число модулей,
    число уроков, наличие обязательных элементов;
  \item \textbf{RAG-метрики}: достоверность относительно источника
    (faithfulness), релевантность ответа (answer relevancy),
    точность контекста (context precision);
  \item \textbf{педагогические метрики}: покрытие уровней таксономии
    Блума, согласованность учебных целей с содержанием курса.
\end{itemize}

\subsection{Идентификация исследовательской лакуны}

Обзор выявил четыре пробела в литературе и практике.

\begin{itemize}
  \item Работы по применению LLM в~подготовке образовательных материалов в~основном
    оценивают связность и полезность текста, но редко связывают
    архитектуру системы с измеримой педагогической метрикой вроде
    покрытия уровней таксономии Блума.
  \item Большинство промышленных ИИ-инструментов для разработки курсов
    ускоряют отдельные этапы работы разработчика курса, но не исследуют мультиагентный
    конвейер с сохраняемым состоянием, который ведёт документ через полный цикл
    методической упаковки и поддерживает человека на ключевых точках принятия решения.
  \item Русскоязычный B2B-контекст представлен слабо: не хватает
    воспроизводимых корпусов, включающих регламенты, продуктово-редакционные
    документы и учебные программы в единой постановке задачи.
  \item Практически не исследована задача автоматического обновления курса
    при изменении исходного документа, хотя именно она определяет
    экономическую ценность ИИ-методиста для корпоративного заказчика.
\end{itemize}

Новизна работы --- не в самом факте использования LLM,
а в экспериментальном исследовании архитектуры, сочетающей RAG,
мультиагентную оркестрацию, педагогическую вариативность
и пригодность к развёртыванию во~внутреннем контуре.

% ─────────────────────────────────────────────────────────────────
\section{Постановка задачи и гипотезы}
% ─────────────────────────────────────────────────────────────────

\subsection{Задача исследования}

Задача: \textit{разработать интеллектуального
методического ассистента, который по~входному документу организации
формирует педагогически состоятельную структуру курса и~черновой контент,
опираясь на~мультиагентную архитектуру с~сохраняемым состоянием
и~обеспечивая приемлемое качество, скорость и~стоимость подготовки.}

\subsection{Методологическая основа}

Постановка опирается на три основания.

\begin{itemize}
  \item \textbf{Working Backwards / CVM}~\cite{amazon_cvm_2014}:
    продуктовые требования формулируются от боли и ценности клиента
    назад к технической реализации.
  \item \textbf{Jobs To Be Done}~\cite{christensen_jtbd_2016}:
    пользователь «нанимает» систему для снятия рутинной нагрузки и~ускорения упаковки знания
    в~обучающий артефакт.
  \item \textbf{Таксономия Блума}~\cite{bloom_taxonomy_1956}:
    качество учебных целей оценивается по глубине когнитивного
    действия, а не только по формальной структуре ответа.
\end{itemize}

\subsection{Три конвейера ИИ-методиста}

Система работает в трёх режимах:

\begin{enumerate}
  \item \textbf{Адаптационный курс}: входом служит продуктовое описание,
    бриф или внутреннее руководство; ключевая задача --- быстро упаковать
    материал для новых сотрудников.
  \item \textbf{Асинхронный учебный курс}: входом являются программа
    курса, учебные материалы или экспертный конспект; здесь критична
    логика модулей, последовательность освоения материала и покрытие учебных целей.
  \item \textbf{Нормативный курс}: входом выступает регламент,
    политика или процедура; здесь на первый план выходит верность
    источнику, работа с фактами и возможность последующего обновления.
\end{enumerate}

\subsection{Исследовательские гипотезы}

Проверяемые гипотезы:

\textbf{H1 (продуктовая).} Мультиагентный конвейер позволяет сократить
время до первого методического черновика до уровня менее 30~минут
на один курс
при сохранении покрытия уровней таксономии Блума не ниже 70\,\%.

\textbf{H2 (техническая).} Мультиагентная архитектура с сохраняемым состоянием
обеспечивает более высокое качество результата по метрикам
faithfulness и покрытия уровней таксономии Блума по сравнению
с одношаговой базовой конфигурацией
при сопоставимой базовой LLM.

\textbf{H3 (бизнесовая).} Себестоимость подготовки одного курса
не превышает 1\,\% от стоимости ручной разработки силами внешнего
подрядчика.

\subsection{Метрики проверки гипотез}

Группы метрик:

\begin{itemize}
  \item \textbf{для H1}: время до первого черновика, доля курсов
    с покрытием уровней таксономии Блума
    $\geq 70\,\%$, доля структурно валидных результатов;
  \item \textbf{для H2}: faithfulness, answer relevancy,
    context precision, покрытие уровней таксономии Блума,
    плотность ссылок на источник;
  \item \textbf{для H3}: себестоимость на курс, отношение себестоимости
    к средней рыночной стоимости подрядной разработки, а также
    параметры подписочной экономики (LTV/CAC, срок окупаемости).
\end{itemize}

Метрики связывают научную часть работы
с практической ценностью: от корректности по отношению
к источнику до экономической оправданности внедрения.
